\noindent The binary heap is faster on every instance, by a factor of \val{abl.DP.sparse.ratio} for the DP on sparse graphs and
\val{abl.G2.sparse.ratio} for G2. This does not contradict \cref{thm:sequence}: the theorem bounds the number of pointer
operations, not the time each one takes. The binary heap keeps its keys in one array that stays in cache, while every Fibonacci
node is a separate allocation with four pointers. On these graphs the asymptotic advantage has little to work with, because a
Dijkstra run on a graph with \(m/n<10\) does fewer decrease-keys than extract-mins (\cref{tab:heapcounts}), and each binary-heap
decrease-key moves a key only a few levels.

The gap does narrow where the theory says it should. On the dense graphs the DP ratio falls to \val{abl.DP.dense.ratio}, and
on \texttt{c17} (\(m/n=25\)) the run makes \val{abl.dp.c17.decr} decrease-keys against \val{abl.dp.c17.extr} extract-mins and
takes \val{abl.dp.c17.fib}\,ms with the Fibonacci heap against \val{abl.dp.c17.bin}\,ms with the binary heap. We report the
Fibonacci heap as the structure that gives the proved bounds, not as the faster choice on SteinLib-sized inputs.
